\documentclass[11pt]{article}
\usepackage{fontspec}
\setmainfont{Times New Roman}
\setsansfont{Arial}
\setmonofont{Consolas}
\usepackage{amsmath,amssymb,mathtools}
\usepackage{geometry}
\usepackage{microtype}
\geometry{a4paper,margin=2.35cm}
\setlength{\parindent}{1.25em}
\setlength{\parskip}{0pt}
\emergencystretch=25em
\sloppy
\hbadness=10000
\vbadness=10000
\hfuzz=2pt
\vfuzz=10pt
\raggedbottom

\begin{document}

% Unit: Paper26 v001. Direct Russian translation from the German final-audited source slice.
\section*{26. Абстрактное построение теории идеалов в алгебраическом числовом поле}

\begin{center}
J. Ber. d. DMV 33 (1924), с. 102
\end{center}

4. E. Noether, Гёттинген: Абстрактное построение теории идеалов в алгебраическом числовом поле.

Были указаны абстрактные свойства, вполне эквивалентные теории идеалов в алгебраическом числовом поле и тем самым характеризующие все кольца, обладающие такой теорией идеалов. Это кольца с единицей и без делителей нуля, в которых выполняется условие двойной цепи: каждая цепь идеалов по делимости обрывается за конечное число шагов, и так же каждая цепь кратных по модулю каждого идеала, отличного от нулевого идеала; кроме того, эти кольца интегрально замкнуты относительно своего поля частных. Если отбросить последнее предположение, то получаются утверждения о теории идеалов в конечном порядке, которые отчасти уже встречаются без доказательства у Dedekind (3-е изд.).

Подробное изложение появится в Mathematische Annalen.

\end{document}
